\documentclass[12pt,a4paper]{article}

\usepackage{amsmath, amssymb, amsthm}
\usepackage{geometry}
\usepackage{hyperref}
\usepackage{graphicx}
\usepackage{booktabs}
\usepackage{cite}
\usepackage{algorithm}
\usepackage{algpseudocode}
\usepackage{xcolor}
\usepackage{microtype}
\usepackage{authblk}
\usepackage{abstract}
\usepackage{titlesec}
\usepackage{enumitem}

\geometry{margin=2.5cm}

\newcommand{\placeholder}[1]{%
  \textcolor{gray}{\textit{[Placeholder: #1]}}%
}

\hypersetup{
  colorlinks=true,
  linkcolor=blue,
  citecolor=blue,
  urlcolor=blue
}

\newtheorem{definition}{Definition}
\newtheorem{proposition}{Proposition}
\newtheorem{remark}{Remark}

\title{\textbf{Spectral Composition: \\
A New Paradigm for\\
Sample-Based Music Creation\\
via Latent Diffusion}}

\author{Lorenzo Moriondo\\
\small{tunedconsulting@gmail.com}\\
\small{lorenzo@genefold.ai}}
\affil{
\large{tuned.org.uk}\\
\large{Genefold AI Ltd}}

\date{\today\\[6pt]
\small\url{https://github.com/tuned-org-uk/latent-sound-diffusion}}

\begin{document}

\maketitle

\begin{abstract}
We introduce \emph{Spectral Composition}, a novel method for sample-based
music production that treats composition as iterative navigation of a
compressed spectral latent manifold rather than direct waveform or
symbolic editing. A frozen ArrowSpace graph Laplacian $L_F$ and its
associated energy-dispersion network $\lambda^{\mathrm{ED}}$, estimated
once from the producer's entire sample library, define a low-dimensional
feature-space manifold over which all generative and compositional
operations are performed. The method proceeds in four stages:
\textbf{(1)~spectral encoding} (``dehydration'') of the library into a
compact spectral representation $(L_F, U_q, \lambda^{\mathrm{ED}})$;
\textbf{(2)~dense unconditional sampling} of a sound bank via a
spectrally-conditioned 1-D diffusion transformer (DiT) operating on
EnCodec audio latents; \textbf{(3)~MIDI-conditioned materialisation}
(``rehydration'') that pitch-shifts and re-timbres generated sounds
according to a language-model-generated MIDI file; and \textbf{(4)~recursive
variant generation} by feeding outputs back through the model with
successive MIDI sequences. The resulting family of variants, each sharing
spectral lineage with the source library, forms a natural set of stems
for multitrack mixing. Spectral Composition is built on the
Latent Sound Diffusion (LSD) system~\cite{lsd}, which is a
sound-production fork of ArrowSpace Latent Diffusion with Spectral Chart
Conditioning (ALD-SC)~\cite{aldsc}, and is exposed to the producer via
the LSD-Studio desktop application~\cite{lsdstudio}.
\end{abstract}

\tableofcontents
\newpage


%% ─────────────────────────────────────────────────────────────────────────────
\section{Introduction}
\label{sec:intro}
%% ─────────────────────────────────────────────────────────────────────────────

Traditional sample-based music production treats a sample library as a
fixed, enumerable set of audio assets: the producer browses, triggers,
slices, and layers clips inside a digital audio workstation (DAW), guided
by perception and taste. This workflow has scaled to industrial size
through large commercial sample libraries, but it fundamentally limits
compositional novelty to recombination of existing clips without access
to the generative structure of the collection.

Spectral Composition proposes a different model: the sample library is a
\emph{corpus} from which a frozen spectral prior is estimated once, after
which composition is an act of sampling and recombination inside the
resulting latent manifold. This reframes ``browsing samples'' as
``traversing a spectral graph'', and ``arranging a track'' as
``recursively rehydrating dense latent draws against symbolic control
signals.''

\subsection{Contributions}

\begin{itemize}[leftmargin=*]
  \item We define the Spectral Composition pipeline: a four-stage method
    (dehydration, sampling, materialisation, recursion) that operationalises
    compositional workflows over a spectral latent manifold.
  \item We show how the existing LSD inference contract
    (\texttt{generate\_sound\_bank}, \texttt{condition\_on\_audio},
    \texttt{synthesize\_midi}) directly implements stages~2--4 of the
    pipeline, requiring only the addition of a language-model MIDI
    front-end.
  \item We articulate the mathematical and architectural foundations of the
    method, grounding them in the ALD-SC research programme's central
    claim: decoding on the feature-space manifold $(L_F,
    \lambda^{\mathrm{ED}})$ yields better global semantic coherence under
    compression than decoding on an unconstrained ambient latent.
  \item We describe LSD-Studio as the producer-facing realisation of the
    pipeline, assigning multitrack mixing and workflow management to the
    application layer.
\end{itemize}

\subsection{Paper Organisation}

Section~\ref{sec:background} reviews background in latent diffusion for
audio, ArrowSpace graph priors, and entropic time scheduling.
Section~\ref{sec:method} defines the four stages of Spectral Composition.
Section~\ref{sec:architecture} describes the underlying model architecture.
Section~\ref{sec:implementation} details the implementation and datasets.
Section~\ref{sec:evaluation} presents evaluation methodology.
Section~\ref{sec:discussion} discusses the method as a creative paradigm.
Section~\ref{sec:futurework} outlines future work, and
Section~\ref{sec:conclusion} concludes.


%% ─────────────────────────────────────────────────────────────────────────────
\section{Background}
\label{sec:background}
%% ─────────────────────────────────────────────────────────────────────────────

\subsection{Latent Diffusion for Audio}
\label{subsec:bg_lda}

Latent diffusion models~\cite{rombach2022} move the expensive iterative
denoising process from pixel (or waveform) space into a compressed
continuous latent space produced by a pretrained autoencoder, drastically
reducing computation while retaining high-fidelity synthesis.
EnCodec~\cite{defossez2022} provides the audio equivalent: a neural codec
operating at 24\,kHz that encodes mono audio clips into a 1-D continuous
latent $z \in \mathbb{R}^{D \times T}$ (with $D=128$, $T=375$ for 5-second
clips) via its pre-quantization continuous features.

\placeholder{Discuss related audio generative models (AudioLM, MusicGen,
Stable Audio, AudioCraft). Contrast with symbolic-only (MIDI) and
waveform-only generation. Motivate the spectral-latent approach.}

\subsection{ArrowSpace Graph Priors and Energy-Dispersion Networks}
\label{subsec:bg_arrowspace}

ArrowSpace~\cite{arrowspace_joss} defines a spectral search framework for
embeddings in which feature columns of a corpus feature matrix
$A \in \mathbb{R}^{N \times F}$ are treated as nodes of a feature-space
graph $G_F = (V_F, W_F)$ with $|V_F|=F$. The feature-space graph
Laplacian
\begin{equation}
  L_F = D^{-1/2}(D - W)D^{-1/2}, \qquad L_F = U\Lambda U^\top,
  \label{eq:laplacian}
\end{equation}
has eigenvectors $U_q \in \mathbb{R}^{F \times q}$ (the $q$ smoothest
modes) that define the semantic subspace of the corpus.

The Energy Dispersion Network (EDN)~\cite{edn2025} associates to each
feature $f$ a per-feature energy-dispersion scalar
$\lambda^{\mathrm{ED}}_f \in [0,1]$, encoding how semantic structure is
distributed over the feature graph. Projected onto the chart, these give
per-mode dispersion weights
\begin{equation}
  \lambda^{\mathrm{chart}}_k
  = \sum_{f=1}^{F} \lambda^{\mathrm{ED}}_f \, U_{f,k}^2.
  \label{eq:chart_dispersion}
\end{equation}

\placeholder{Expand on the ArrowSpace spectral search paper, the kNN
graph construction, $\tau$-modulation, and the connection between
$\lambda^{\mathrm{ED}}$ and semantic coherence under compression.}

\subsection{Entropic Time Scheduling and the Barontini Clock}
\label{subsec:bg_clock}

The entropic (Barontini) clock~\cite{barontini2026,stancevic2025} provides
an intrinsic, spectrally-informed stopping criterion for diffusion
sampling. Per-mode entropy-time coordinates are defined as
\begin{equation}
  \tau_k(t) = \nu_k \bar{\alpha}_k(t),
  \label{eq:entropic_time}
\end{equation}
where $\nu_k$ are the eigenvalues of $L_F$ and $\bar{\alpha}_k(t)$ is the
mode-$k$ noise schedule. The sampler terminates when
\begin{equation}
  \sum_{k=1}^{q} \nu_k \, \bar{\alpha}_k(t) < \varepsilon,
  \label{eq:heat_death}
\end{equation}
the ``heat-death'' criterion, replacing an arbitrary fixed step count
with a spectrally-motivated halting rule~\cite{lsd,aldsc}.

\placeholder{Discuss the connection to the problem of time in quantum
gravity and cold-atom experiments (Barontini 2026). Relate to prior work
on adaptive diffusion schedules (DDIM, DPM-Solver, etc.).}


%% ─────────────────────────────────────────────────────────────────────────────
\section{Spectral Composition}
\label{sec:method}
%% ─────────────────────────────────────────────────────────────────────────────

We now define the four stages of the Spectral Composition pipeline
formally. Let $\mathcal{X} = \{x_1, \ldots, x_N\}$ denote the producer's
sound archive (a set of raw audio clips), and let $\mathcal{M}$ denote a
trained LSD model (frozen prior + graph decoder + 1-D DiT).

\subsection{Stage 1 — Spectral Encoding (Library Dehydration)}
\label{subsec:stage1}

\subsubsection{EnCodec feature extraction}

Each clip $x_i \in \mathcal{X}$ is encoded by the frozen EnCodec encoder
$E_\phi$ (24\,kHz mono, never updated) to produce:
\begin{equation}
  (z_i,\; A_i) = E_\phi(x_i),
  \qquad
  z_i \in \mathbb{R}^{D \times T},\quad
  A_i = \mathrm{Pool}(z_i) \in \mathbb{R}^F.
  \label{eq:encodec}
\end{equation}
The pooled feature field $A_i$ is the row contributed by clip $i$ to the
corpus feature matrix $A \in \mathbb{R}^{N \times F}$.

\subsubsection{ArrowSpace prior construction}

From $A$ the ArrowSpace adapter builds the frozen feature-space graph
Laplacian $L_F$ and energy-dispersion network $\lambda^{\mathrm{ED}}$
via the \texttt{pyarrowspace} library~\cite{pyarrowspace}. The eigenvectors
$U_q$ and eigenvalues $\nu_k$ are stored as non-trainable buffers. The
spectral conditioning vector for a clip is
\begin{equation}
  c_{\mathrm{spec}} =
  \bigl[\tilde{e},\; \lambda^{\mathrm{chart}},\; \nu\bigr]
  \in \mathbb{R}^{3q},
  \label{eq:cspec}
\end{equation}
where $\tilde{e}_k$ is the per-mode band energy. The entire library is
now represented by the compact frozen objects
$(L_F, U_q, \lambda^{\mathrm{ED}})$ — the \emph{dehydrated} spectral
representation.

\subsubsection{Two-stage model training}

Training proceeds as two frozen-then-fine-tuned phases.
In Phase 1 the graph decoder is trained with optional latent-space noise
injection (parameter \texttt{NOISE\_INJECT} $\in [0, \infty)$).
In Phase 2 the 1-D DiT denoiser is trained on the corpus latents
$\{z_i\}$ conditioned on $c_{\mathrm{spec}}$:
\begin{equation}
  \mathcal{L}_{\mathrm{diff}}
  = \mathbb{E}_{z_0,\epsilon,t}
    \left\|v - v_\theta(z_t,\, t,\, c_{\mathrm{spec}})\right\|_2^2,
  \label{eq:diffusion_loss}
\end{equation}
using $v$-prediction for numerical stability.

\placeholder{Describe the full training loss including $\mathcal{L}_{\mathrm{rec}}$
(L1 + multi-scale STFT), $\mathcal{L}_{\mathrm{chart}}$, and
$\mathcal{L}_{\mathrm{smooth}}$. Discuss the non-repeatable training
design (SEED=None, NOISE\_INJECT).}

\subsection{Stage 2 — Dense Unconditional Sampling (Sound-Bank Generation)}
\label{subsec:stage2}

Given trained $\mathcal{M}$, a sound bank $\mathcal{B}$ of $n$ clips is
generated by running the DDIM or Euler sampler on the 1-D DiT
unconditionally, then decoding through the
\texttt{ClockGatedGraphDecoder}:
\begin{equation}
  \hat{z}^{(j)} \sim p_\theta(z \mid c_{\mathrm{spec}}),
  \qquad j = 1, \ldots, n,
  \label{eq:sampling}
\end{equation}
\begin{equation}
  \hat{x}^{(j)} = D_\psi\!\left(\hat{z}^{(j)};\;
    U_q,\; \lambda^{\mathrm{ED}},\; c_{\mathrm{spec}}^{(j)}\right),
  \label{eq:decoding}
\end{equation}
where $c_{\mathrm{spec}}^{(j)}$ is derived \emph{self-consistently} from
$\hat{z}^{(j)}$ itself (no separate chart sampling required). The bank
$\mathcal{B} = \{\hat{x}^{(1)}, \ldots, \hat{x}^{(n)}\}$ consists of
samples that are perceptually self-similar — drawn from the same frozen
manifold — yet acoustically distinct because \texttt{SEED=None} is the
default.

\placeholder{Discuss the ``super-dense'' character of the bank: why
all samples share stylistic identity while varying locally. Relate to the
frozen prior design constraint (corpus-level prior, not per-clip).}

\subsection{Stage 3 — MIDI-Conditioned Materialisation (Rehydration)}
\label{subsec:stage3}

\subsubsection{Language-model MIDI generation}

A language model (e.g.\ a GPT-based symbolic music model) generates a MIDI
file $\mathcal{P}$ = $\{(p_i, t_i, d_i)\}$ where $p_i$ is the MIDI pitch,
$t_i$ the onset time, and $d_i$ the duration of note $i$.

\placeholder{Describe the GPT-MIDI generation front-end: prompt design,
conditioning on genre/key/BPM, and the interface to LSD-Studio. Discuss
alternative symbolic sources (human MIDI, rule-based generation, MusicXML).}

\subsubsection{Pitch-shifted decoding (synthesize\_midi)}

For each note $(p_i, t_i, d_i)$ the adapter selects a bank clip
$\hat{x}^{(j)}$ and applies pitch-shifting to MIDI pitch $p_i$ via the
\texttt{synthesize\_midi} method of \texttt{LSDModel}:
\begin{equation}
  y_i = \mathrm{PitchShift}\!\left(\hat{x}^{(j)},\;
    \Delta p_i\right)\bigg|_{[t_i,\; t_i + d_i]},
  \label{eq:pitchshift}
\end{equation}
where $\Delta p_i$ is the semitone offset from the bank clip's reference
pitch. The materialised output is
\begin{equation}
  Y^{(1)} = \sum_i y_i,
  \label{eq:materialise}
\end{equation}
a waveform that follows the symbolic structure of $\mathcal{P}$ while
carrying the spectral identity of the source library.

\subsubsection{Audio-conditioned variant generation}

Optionally, the materialised output $Y^{(1)}$ is fed back into the model
via \texttt{condition\_on\_audio} (Mode~B) to generate variants:
\begin{equation}
  \hat{z}' \sim p_\theta(z \mid z_{\mathrm{cond}} + \eta),
  \qquad \eta \sim \mathcal{N}(0, \sigma^2 I),
  \label{eq:conditioned_sampling}
\end{equation}
where $z_{\mathrm{cond}} = E_\phi(Y^{(1)})$ is the conditioning latent
and $\sigma$ is the perturbation strength. The chart conditioning
$c_{\mathrm{spec}}$ is re-derived from $\hat{z}'$, keeping the output on
the archive manifold.

\subsection{Stage 4 — Recursive Variant Generation}
\label{subsec:stage4}

Stages~2 and~3 are applied recursively. Let $\mathcal{P}^{(1)}, \ldots,
\mathcal{P}^{(R)}$ be $R$ distinct MIDI files (generated by the language
model or supplied by the producer). The recursive scheme is:
\begin{align}
  \mathcal{B}^{(0)} &= \texttt{generate\_sound\_bank}(n),
  \label{eq:recursive_init}\\
  Y^{(r)} &= \texttt{synthesize\_midi}\!\left(\mathcal{P}^{(r)},\;
    \mathcal{B}^{(r-1)}\right),
  \label{eq:recursive_synth}\\
  \mathcal{B}^{(r)} &= \texttt{condition\_on\_audio}\!\left(Y^{(r)},\;
    n,\; \sigma_r\right),
  \qquad r = 1, \ldots, R.
  \label{eq:recursive_cond}
\end{align}
Each round deepens the stylistic imprint of the source library while
incorporating new symbolic content from $\mathcal{P}^{(r)}$. The family
$\{Y^{(1)}, \ldots, Y^{(R)}\}$ constitutes a set of \emph{variants}
related by shared spectral ancestry.

\subsubsection{Convergence and heat-death stopping}

The recursive process is bounded intrinsically by the heat-death criterion
(Eq.~\ref{eq:heat_death}): if the conditioning latent $z_{\mathrm{cond}}$
has collapsed to a near-constant spectral chart, additional conditioning
passes yield diminishing novelty. In practice, the producer controls depth
via $R$ and the perturbation schedule $\{\sigma_r\}$.

\subsection{Final Assembly — Multitrack Mixing}
\label{subsec:mixing}

The variant outputs $\{Y^{(1)}, \ldots, Y^{(R)}\}$ plus any selections
from the unconditional bank $\mathcal{B}^{(0)}$ form a natural stem set
for a multitrack mix. Within LSD-Studio, these are routed to a mix bus
with per-stem gain, panning, and effect chains. Additional post-processing
(overlap-and-add long-form generation, BPS-synchronised modulation,
latent-space concatenation) is available via the notebook series
(02--04)~\cite{lsd}.

\placeholder{Describe the DAW integration and MIDI synchronisation in
LSD-Studio. Discuss how BPM and key are propagated from the LM-generated
MIDI through to the mix.}


%% ─────────────────────────────────────────────────────────────────────────────
\section{System Architecture}
\label{sec:architecture}
%% ─────────────────────────────────────────────────────────────────────────────

\subsection{Frozen ArrowSpace Prior and Manifold Construction}
\label{subsec:arch_prior}

The ArrowSpace prior is built once from the corpus by \texttt{build\_prior.py}
and stored as non-trainable buffers:
\begin{equation}
  \Pi_q = U_q U_q^\top \in \mathbb{R}^{F \times F},
  \label{eq:projector}
\end{equation}
the spectral projector onto the smooth subspace. The frozen design
constraint means that $L_F$ and $U_q$ never receive gradient updates;
only the graph decoder and DiT are trained. This ensures that the
manifold geometry is a property of the \emph{corpus}, not a function of
any single training run.

\placeholder{Describe the kNN graph construction, the ArrowSpace library
API (\texttt{gl}, \texttt{lambdas()}), and the fallback behaviour when
\texttt{pyarrowspace} bindings are unavailable.}

\subsection{1-D DiT Denoiser}
\label{subsec:arch_dit}

The denoiser \texttt{MinimalDiT} patchifies the 1-D latent
$z \in \mathbb{R}^{D \times T}$ with stride-$s$ convolutions, applies
$L$ transformer blocks with AdaLN conditioning on $(t, c_{\mathrm{spec}})$,
and unpatchifies to predict velocity $v$. Classifier-free guidance (CFG)
is applied by randomly dropping $c_{\mathrm{spec}}$ during training with
probability $p_{\mathrm{drop}}$, enabling guidance-scale control at
inference.

\subsection{Graph Decoder}
\label{subsec:arch_decoder}

The \texttt{ClockGatedGraphDecoder} processes the denoised latent
$\hat{z}$ through a sequence of \texttt{WaveReconstructionBlock}s:
\begin{align}
  \hat{H} &= A \cdot U_q,
  \label{eq:project_chart}\\
  g &= \sigma(W_g \cdot c_{\mathrm{spec}})\; \odot\; \bar{\alpha}(t),
  \label{eq:gate}\\
  A' &= (\hat{H} \odot g) \cdot U_q^\top + A_{\mathrm{res}},
  \label{eq:lift}
\end{align}
where $A_{\mathrm{res}}$ is a residual convolution path and
$\bar{\alpha}(t)$ is the Barontini clock gate (weak early in denoising,
strong late). The decoder thus allocates reconstruction energy according
to the corpus dispersion $\lambda^{\mathrm{ED}}$, not arbitrary learned
weights.

\subsection{Inference Contract (A / B / C)}
\label{subsec:arch_inference}

\texttt{LSDModel} (defined in \texttt{src/ald\_sc/inference.py})
exposes three methods that directly implement Stages~2--4:

\begin{description}
  \item[Mode A: \texttt{generate\_sound\_bank(n, steps, temperature, seed)}]
    Unconditional DiT sampling + graph decode. Implements Stage~2.
  \item[Mode B: \texttt{condition\_on\_audio(audio, n, strength, seed)}]
    Conditioned sampling with SDE bridge toward the input latent.
    Used in Stage~4 feedback loops.
  \item[Mode C: \texttt{synthesize\_midi(midi\_path, bank, seed)}]
    Pitch-shifted bank decoding driven by a MIDI file.
    Implements Stage~3.
\end{description}

\subsection{LSD-Studio: Producer UX and Mix Bus}
\label{subsec:arch_studio}

LSD-Studio~\cite{lsdstudio} is a Rust + \texttt{iced} desktop
application that wraps the Python LSD model workers and exposes the
full Spectral Composition workflow to the producer: library ingestion,
model baking, sound-bank auditioning, MIDI sequencing, recursive
generation, and multitrack mixing. The application builds to a single
self-contained binary; Python workers are embedded at compile time.

\placeholder{Architecture diagram of the IPC bridge between Rust frontend
and Python workers. Describe the preset system and mix-bus signal flow.}


%% ─────────────────────────────────────────────────────────────────────────────
\section{Implementation and Datasets}
\label{sec:implementation}
%% ─────────────────────────────────────────────────────────────────────────────

\subsection{Software Stack}
\label{subsec:impl_stack}

The LSD model core is implemented in Python~3.13 with PyTorch~$\geq$2.2,
using \texttt{uv} for dependency management. The graph prior is computed
via \texttt{pyarrowspace} (Rust bindings). LSD-Studio is implemented in
Rust (stable) with the \texttt{iced} GUI framework and communicates with
Python workers over a JSON-RPC IPC channel. The full source code,
evaluation pipeline, and trained-model checkpoints are available at
\url{https://github.com/tuned-org-uk/latent-sound-diffusion}.

\subsection{Sound Archive and Licensing}
\label{subsec:impl_datasets}

LSD-Studio ships with support for four commercially-licensed datasets,
all carrying CC0 or CC~BY~4.0 licences suitable for commercial
production:

\begin{table}[h]
\centering
\caption{Supported sound archive datasets.}
\begin{tabular}{llll}
\toprule
Dataset & Size & Content & Licence \\
\midrule
NSynth~\cite{nsynth} & 300K notes & Musical notes, 1006 instruments & CC BY 4.0 \\
Clotho~\cite{clotho} & $\sim$5K clips & Environmental audio with captions & CC BY 4.0 \\
FSD50K-CC0 & $\sim$10K clips & Diverse environmental audio & CC0 \\
Freesound Loop Dataset & $\sim$9700 clips & Loop clips & CC0 \\
\bottomrule
\end{tabular}
\label{tab:datasets}
\end{table}

\subsection{Training Details}
\label{subsec:impl_training}

Two design decisions make each training run artistically non-repeatable:
\textbf{(i)}~\texttt{SEED=None}: a fresh random seed is sampled at the
start of each run, so no two trained models are identical;
\textbf{(ii)}~\texttt{NOISE\_INJECT}: a scalar $\sigma_{\mathrm{inj}} \geq 0$
controls noise injected into the EnCodec latent $z$ before the graph
decoder's training pass, increasing acoustic variety at the cost of
reconstruction fidelity.

\placeholder{Table of hyperparameters (learning rate, batch size,
decoder/DiT architecture dimensions, $q$, $K$, \texttt{NOISE\_INJECT}
sweep). Training time on target hardware.}


%% ─────────────────────────────────────────────────────────────────────────────
\section{Evaluation}
\label{sec:evaluation}
%% ─────────────────────────────────────────────────────────────────────────────

\subsection{Perceptual Coherence of Recursive Variants}
\label{subsec:eval_recursive}

\placeholder{
Propose a listening study evaluating whether recursive variants are
perceived as stylistically related to the source library across $R$
rounds. Metrics: pairwise timbral similarity (e.g.\ CLAP embedding
cosine distance), spectral centroid/rolloff drift over rounds, producer
preference ratings.
}

\subsection{Reconstruction Fidelity: Graph Decoder vs.\ Baseline}
\label{subsec:eval_reconstruction}

The primary controlled comparison isolates the graph decoder
(\texttt{WaveReconstructionBlock} + $U_q$ + $\lambda^{\mathrm{ED}}$) against
a baseline decoder of identical channel widths and upsampling strides but
with plain \texttt{ResBlock1d} (no $U_q$, no $\lambda^{\mathrm{ED}}$).

\placeholder{
Report multi-scale STFT reconstruction loss, audio FAD (once dependency
issues with \texttt{fadtk} are resolved), and per-band spectral energy
retention for graph vs.\ baseline on ESC-50 and NSynth subsets.
}

\subsection{Ablation Studies}
\label{subsec:eval_ablation}

\placeholder{
Ablation over: (1) $q$ (number of retained spectral modes), (2)
\texttt{NOISE\_INJECT} level, (3) CFG guidance scale, (4) heat-death
threshold $\varepsilon$, (5) recursion depth $R$. Metrics: diversity
(mean pairwise distance in CLAP space) vs.\ fidelity (STFT loss).
}

\subsection{Limitations and Open Issues}
\label{subsec:eval_limitations}

\placeholder{
Discuss: FAD computation (\texttt{fadtk} removed; multi-scale STFT used
as fallback); real-data experiments pending (ESC-50, NSynth);
pitch-shifting artefacts in \texttt{synthesize\_midi}; scalability of
ArrowSpace graph construction to large archives ($N > 10$K); and the
absence of long-form ($>$8s) compositional structure in the current
recursive scheme.
}


%% ─────────────────────────────────────────────────────────────────────────────
\section{Discussion}
\label{sec:discussion}
%% ─────────────────────────────────────────────────────────────────────────────

\subsection{Spectral Composition as a Creative Paradigm}
\label{subsec:disc_paradigm}

\placeholder{
Contrast Spectral Composition with (a) loop-based composition (static
library, manual arrangement), (b) sample-based AI tools that generate
independent clips (AudioCraft, Stable Audio), and (c) purely symbolic
generative models (MusicGen, Jukebox). Argue that the manifold-constrained,
recursively rehydrated workflow occupies a distinct position: it maintains
strong stylistic identity with the producer's own archive while enabling
genuine novelty.
}

\subsection{Relationship to Generative-AI Music Tools}
\label{subsec:disc_related}

\placeholder{
Survey of relevant systems: Google MusicLM, Meta AudioCraft/MusicGen,
Stability Audio, Riffusion, Moûsai, AudioLDM 2. Situate Spectral
Composition relative to each with respect to: (1) source identity
preservation, (2) producer control, (3) symbolic grounding via MIDI,
(4) recursive/iterative generation.
}

\subsection{Ethical and Licensing Considerations}
\label{subsec:disc_ethics}

\placeholder{
Discuss: training on commercially-licensed datasets (all CC0/CC~BY~4.0);
the dehydration step as a form of model-mediated compression of the
producer's private library; implications for sample clearance; and the
``sound of the future'' framing as a creative rather than extractive
application.
}


%% ─────────────────────────────────────────────────────────────────────────────
\section{Future Work}
\label{sec:futurework}
%% ─────────────────────────────────────────────────────────────────────────────

\subsection{Text/CLAP Conditioning and Genre Control}

\placeholder{
Extend Stage~2 and Stage~3 with text and CLAP conditioning (tracked as
LSD Phase~2). Discuss joint text + graph CFG guidance and genre-conditioned
bank generation.
}

\subsection{Long-Form Generation and Wave Recurrence}

\placeholder{
Describe the overlap-and-add notebook (02) and latent-space concatenation
(04) as current approximations to long-form composition. Motivate the
ESDM wave-recurrence mechanism~\cite{esdm} as a principled replacement,
enabling second-order graph-wave propagation in the decoder.
}

\subsection{Full Entropic Training Schedules}

\placeholder{
The current model uses the entropic clock only at inference (sampling
stopping criterion). Discuss training with per-mode noise schedules as
an open research direction~\cite{aldsc} (issue \#11).
}

\subsection{Federated and Personalised Priors}

\placeholder{
Motivate federated training of ArrowSpace priors across multiple producers'
private archives without sharing raw audio, as a privacy-preserving
extension of the dehydration step.
}


%% ─────────────────────────────────────────────────────────────────────────────
\section{Conclusion}
\label{sec:conclusion}
%% ─────────────────────────────────────────────────────────────────────────────

Spectral Composition introduces a four-stage compositional workflow —
dehydration, dense sampling, MIDI-conditioned materialisation, and
recursive variant generation — that grounds sample-based music production
in a compressed spectral latent manifold estimated from the producer's
own library. The method is fully implemented in the LSD model core and
LSD-Studio application, requiring only the addition of a language-model
MIDI front-end to complete the pipeline. By anchoring every generative
and compositional operation to the same frozen ArrowSpace prior
$(L_F, U_q, \lambda^{\mathrm{ED}})$, Spectral Composition ensures that
all generated and rehydrated sounds share spectral lineage with the
source archive, yielding variant families that are simultaneously novel
and stylistically coherent.

\placeholder{Summarise the main experimental findings (to be filled once
evaluation is complete).}


%% ─────────────────────────────────────────────────────────────────────────────
\bibliographystyle{plain}
\begin{thebibliography}{99}

\bibitem{lsd}
tuned-org-uk.
\newblock {Latent Sound Diffusion (LSD)}: Sound-generation fork of ALD-SC.
\newblock GitHub repository, 2026.
\newblock \url{https://github.com/tuned-org-uk/latent-sound-diffusion}

\bibitem{aldsc}
tuned-org-uk.
\newblock {ArrowSpace Latent Diffusion with Spectral Chart Conditioning (ALD-SC)}.
\newblock GitHub repository, 2026.
\newblock \url{https://github.com/tuned-org-uk/arrowspace-latent-diffusion}

\bibitem{lsdstudio}
tuned-org-uk.
\newblock {LSD Studio}: Latent Sound Diffusion desktop application.
\newblock GitHub repository, 2026.
\newblock \url{https://github.com/tuned-org-uk/LSD-studio}

\bibitem{esdm}
tuned-org-uk.
\newblock {Entropic Semantic Diffusion Model (ESDM)}.
\newblock GitHub repository, 2026.
\newblock \url{https://github.com/tuned-org-uk/entropic-semantic-diffusion}

\bibitem{pyarrowspace}
tuned-org-uk.
\newblock {pyarrowspace}: ArrowSpace library (Rust bindings).
\newblock GitHub repository, 2026.
\newblock \url{https://github.com/tuned-org-uk/pyarrowspace}

\bibitem{arrowspace_joss}
L.\ Moriondo et al.
\newblock {ArrowSpace}: Spectral search for embeddings.
\newblock \textit{Journal of Open Source Software}, 2026.
\newblock \url{https://doi.org/10.21105/joss.09002}

\bibitem{edn2025}
tuned-org-uk.
\newblock Energy Dispersion Networks.
\newblock arXiv:2606.21535, 2025.
\newblock \url{https://arxiv.org/abs/2606.21535}

\bibitem{rombach2022}
R.\ Rombach, A.\ Blattmann, D.\ Lorenz, P.\ Esser, and B.\ Ommer.
\newblock High-resolution image synthesis with latent diffusion models.
\newblock In \textit{CVPR}, 2022.

\bibitem{defossez2022}
A.\ D\'{e}fossez, J.\ Copet, G.\ Synnaeve, and Y.\ Adi.
\newblock {EnCodec}: High fidelity neural audio compression.
\newblock arXiv:2210.13438, 2022.

\bibitem{barontini2026}
F.\ Barontini.
\newblock Testing the problem of time with cold atoms.
\newblock \textit{Physical Review Letters}, 2026.

\bibitem{stancevic2025}
I.\ Stancevic et al.
\newblock Entropic time schedulers for generative diffusion models.
\newblock arXiv preprint, 2025.

\bibitem{nsynth}
J.\ Engel et al.
\newblock Neural audio synthesis of musical notes with {WaveNet} autoencoders.
\newblock In \textit{ICML}, 2017.

\bibitem{clotho}
K.\ Drossos, S.\ Lipping, and T.\ Virtanen.
\newblock Clotho: An audio captioning dataset.
\newblock In \textit{ICASSP}, 2020.

\bibitem{piczak2015}
K.\ J.\ Piczak.
\newblock {ESC}: Dataset for environmental sound classification.
\newblock In \textit{ACM Multimedia}, 2015.

\bibitem{spectral_blog}
tuned.org.uk.
\newblock Diffusion as spectral-geometric projection, 2024.
\newblock \url{https://www.tuned.org.uk/posts/021_diffusion_as_spectral_geometric_projection/}

\end{thebibliography}

\end{document}